\documentclass[12pt]{article}
\usepackage[utf8]{inputenc}
\usepackage[russian]{babel}
\usepackage{graphicx}
\usepackage{lipsum}
\usepackage{tikz}

\begin{document}

\title{Демонстрация работы с изображениями и ссылками в LaTeX}
\author{Демонстрационный документ}
\date{\today}
\maketitle

\tableofcontents

\section{Введение}
\label{sec:intro}
Это демонстрационный документ для изучения различных возможностей LaTeX.

\lipsum[1-2]

\section{Работа с изображениями}
\label{sec:images}

\subsection{Создание собственного изображения с TikZ}
\label{subsec:tikz}

Создадим собственное изображение с помощью TikZ:

\begin{figure}[htbp]
    \centering
    \begin{tikzpicture}
        % Рисуем дом
        \draw[fill=yellow!30] (0,0) rectangle (4,3);
        \draw[fill=red!50] (0,3) -- (2,5) -- (4,3) -- cycle;
        \draw[fill=brown!60] (1.5,0) rectangle (2.5,1.5);
        \draw[fill=cyan!30] (0.5,1.5) rectangle (1.3,2.3);
        \draw[fill=cyan!30] (2.7,1.5) rectangle (3.5,2.3);
        % Солнце
        \draw[fill=yellow] (6,4) circle (0.5);
        \foreach \angle in {0,45,...,315}
            \draw[yellow, thick] (6,4) -- +(\angle:0.8);
    \end{tikzpicture}
    \caption{Собственное изображение, созданное с помощью TikZ}
    \label{fig:myhouse}
\end{figure}

Ссылка на рисунок~\ref{fig:myhouse} на странице~\pageref{fig:myhouse}.

\subsection{Эксперименты с параметрами изображений}
\label{subsec:params}

\subsubsection{Использование height и width}
\label{subsubsec:heightwidth}

\lipsum[3]

\begin{figure}[h]
    \centering
    \begin{tikzpicture}
        \draw[fill=blue!20] (0,0) circle (1);
        \node at (0,0) {width=5cm};
    \end{tikzpicture}
    \caption{Изображение с фиксированной шириной 5cm}
    \label{fig:width5cm}
\end{figure}

\begin{figure}[h]
    \centering
    \begin{tikzpicture}[scale=0.5]
        \draw[fill=green!20] (0,0) circle (1);
        \node at (0,0) {height=3cm};
    \end{tikzpicture}
    \caption{Изображение с фиксированной высотой 3cm}
    \label{fig:height3cm}
\end{figure}

\subsubsection{Использование angle и scale}
\label{subsubsec:anglescale}

\begin{figure}[h]
    \centering
    \begin{tikzpicture}[rotate=45]
        \draw[fill=red!20] (0,0) rectangle (2,1);
        \node at (1,0.5) {angle=45};
    \end{tikzpicture}
    \caption{Изображение повёрнуто на 45 градусов}
    \label{fig:angle45}
\end{figure}

\begin{figure}[h]
    \centering
    \begin{tikzpicture}[scale=2]
        \draw[fill=purple!20] (0,0) rectangle (1,1);
        \node at (0.5,0.5) {scale=2};
    \end{tikzpicture}
    \caption{Изображение увеличено в 2 раза}
    \label{fig:scale2}
\end{figure}

\lipsum[4]

\section{Работа с textwidth и linewidth}
\label{sec:widths}

\subsection{Использование textwidth}
\label{subsec:textwidth}

\begin{figure}[htbp]
    \centering
    \begin{tikzpicture}
        \draw[fill=orange!30] (0,0) rectangle (0.8\textwidth, 2);
        \node at (0.4\textwidth, 1) {width=0.8\textbackslash textwidth};
    \end{tikzpicture}
    \caption{Изображение шириной 0.8 от textwidth}
    \label{fig:textwidth}
\end{figure}

\lipsum[5]

\subsection{Использование linewidth}
\label{subsec:linewidth}

\begin{figure}[htbp]
    \centering
    \begin{tikzpicture}
        \draw[fill=teal!30] (0,0) rectangle (0.6\linewidth, 1.5);
        \node[text width=0.5\linewidth, align=center] at (0.3\linewidth, 0.75) 
            {width=0.6\textbackslash linewidth};
    \end{tikzpicture}
    \caption{Изображение шириной 0.6 от linewidth}
    \label{fig:linewidth}
\end{figure}

\lipsum[6]

\section{Эксперименты с плавающими объектами}
\label{sec:floats}

\lipsum[7]

\begin{figure}[h]
    \centering
    \begin{tikzpicture}
        \draw[fill=pink!30] (0,0) circle (1);
        \node at (0,0) {[h]};
    \end{tikzpicture}
    \caption{Плавающий объект с спецификатором [h] - здесь}
    \label{fig:here}
\end{figure}

\lipsum[8]

\begin{figure}[t]
    \centering
    \begin{tikzpicture}
        \draw[fill=lime!30] (0,0) circle (1);
        \node at (0,0) {[t]};
    \end{tikzpicture}
    \caption{Плавающий объект с спецификатором [t] - вверху страницы}
    \label{fig:top}
\end{figure}

\lipsum[9]

\begin{figure}[b]
    \centering
    \begin{tikzpicture}
        \draw[fill=violet!30] (0,0) circle (1);
        \node at (0,0) {[b]};
    \end{tikzpicture}
    \caption{Плавающий объект с спецификатором [b] - внизу страницы}
    \label{fig:bottom}
\end{figure}

\lipsum[10]

\begin{figure}[p]
    \centering
    \begin{tikzpicture}
        \draw[fill=cyan!30] (0,0) circle (1);
        \node at (0,0) {[p]};
    \end{tikzpicture}
    \caption{Плавающий объект с спецификатором [p] - на отдельной странице}
    \label{fig:page}
\end{figure}

\lipsum[11]

\begin{figure}[htbp]
    \centering
    \begin{tikzpicture}
        \draw[fill=magenta!30] (0,0) circle (1);
        \node at (0,0) {[htbp]};
    \end{tikzpicture}
    \caption{Плавающий объект с спецификатором [htbp] - LaTeX выбирает лучшее место}
    \label{fig:flexible}
\end{figure}

\lipsum[12]

\section{Тестирование команды label}
\label{sec:labeltests}

\subsection{Проверка количества прогонов}
\label{subsec:runs}

Эта секция~\ref{sec:labeltests} содержит подсекцию~\ref{subsec:runs}.

\subsubsection{Первый подподраздел}
\label{subsubsec:first}

\lipsum[13]

\subsubsection{Второй подподраздел}
\label{subsubsec:second}

Нумерованный список:
\begin{enumerate}
    \item Первый элемент \label{item:first}
    \item Второй элемент \label{item:second}
    \item Третий элемент \label{item:third}
\end{enumerate}

Ссылка на элемент~\ref{item:second} списка.

\lipsum[14]

\subsection{Label перед caption}
\label{subsec:wronglabel}

\textbf{ЭКСПЕРИМЕНТ: Label перед caption}

\begin{figure}[htbp]
    \centering
    \label{fig:wronglabel} % НЕПРАВИЛЬНО: label перед caption
    \begin{tikzpicture}
        \draw[fill=red!50] (0,0) rectangle (3,2);
        \node[text width=2.5cm, align=center] at (1.5,1) {Label ПЕРЕД caption};
    \end{tikzpicture}
    \caption{Рисунок с неправильным размещением label}
\end{figure}

\begin{figure}[htbp]
    \centering
    \begin{tikzpicture}
        \draw[fill=green!50] (0,0) rectangle (3,2);
        \node[text width=2.5cm, align=center] at (1.5,1) {Label ПОСЛЕ caption};
    \end{tikzpicture}
    \caption{Рисунок с правильным размещением label}
    \label{fig:correctlabel} % ПРАВИЛЬНО: label после caption
\end{figure}

\textbf{Результат:}
\begin{itemize}
    \item Ссылка на рисунок с label перед caption: \ref{fig:wronglabel} (может быть неправильной!)
    \item Ссылка на рисунок с label после caption: \ref{fig:correctlabel} (правильная)
\end{itemize}

\lipsum[15]

\section{Эксперименты с уравнениями}
\label{sec:equations}

\subsection{Правильное размещение label}
\label{subsec:correcteq}

Правильный способ - label внутри окружения equation:

\begin{equation}
    E = mc^2
    \label{eq:einstein}
\end{equation}

Ссылка на уравнение~\ref{eq:einstein} работает правильно.

\subsection{Неправильное размещение label}
\label{subsec:wrongeq}

\textbf{ЭКСПЕРИМЕНТ: Label после end\{equation\}}

\begin{equation}
    a^2 + b^2 = c^2
\end{equation}
\label{eq:pythagoras} % НЕПРАВИЛЬНО: label после end{equation}

\textbf{Результат:} Ссылка на уравнение~\ref{eq:pythagoras} может указывать на секцию, а не на уравнение!

Правильный вариант:

\begin{equation}
    \int_{-\infty}^{\infty} e^{-x^2} dx = \sqrt{\pi}
    \label{eq:gaussian}
\end{equation}

Ссылка на уравнение~\ref{eq:gaussian} работает корректно.

\lipsum[16]

\section{Заключение}
\label{sec:conclusion}

\subsection{Выводы по экспериментам}
\label{subsec:findings}

\begin{enumerate}
    \item \textbf{Параметры изображений:}
    \begin{itemize}
        \item width и height задают абсолютные размеры
        \item angle поворачивает изображение
        \item scale масштабирует изображение
    \end{itemize}
    
    \item \textbf{textwidth vs linewidth:}
    \begin{itemize}
        \item textwidth - ширина текстового блока
        \item linewidth - ширина текущей строки (меняется в многоколоночном режиме)
    \end{itemize}
    
    \item \textbf{Спецификаторы позиции плавающих объектов:}
    \begin{itemize}
        \item h - здесь (here)
        \item t - вверху страницы (top)
        \item b - внизу страницы (bottom)
        \item p - на отдельной странице (page)
        \item ! - игнорировать некоторые ограничения
    \end{itemize}
    
    \item \textbf{Количество прогонов:} Обычно требуется 2-3 прогона для корректной работы всех ссылок
    
    \item \textbf{Размещение label:}
    \begin{itemize}
        \item Для рисунков и таблиц: label ПОСЛЕ caption
        \item Для уравнений: label ВНУТРИ окружения equation
        \item Неправильное размещение приводит к неверным ссылкам
    \end{itemize}
\end{enumerate}

\lipsum[17]

\end{document}



